\documentclass[11pt]{article}

\usepackage[a4paper,margin=0.88in]{geometry}
\usepackage{amsmath,amssymb,amsthm}
\usepackage{array,booktabs,graphicx}
\usepackage[hidelinks]{hyperref}
\usepackage{microtype}
\usepackage[compact]{titlesec}
\titlespacing*{\section}{0pt}{1.4ex plus .3ex}{.7ex}

\newtheorem{theorem}{Theorem}
\newtheorem{lemma}{Lemma}
\newcommand{\R}{\mathbb R}
\newcommand{\Sym}{\operatorname{Sym}}
\newcommand{\rank}{\operatorname{rank}}
\newcommand{\spn}{\operatorname{span}}

\title{A Clifford Construction Answering Problem 5.9\\
in \emph{Phase Retrieval by Projections}}
\author{Automatic math-discovery run \texttt{fa\_banach\_001}}
\date{August 11, 2026}

\begin{document}
\maketitle

\section*{Source problem}

\noindent\textbf{Source:} \emph{Phase Retrieval by Projections},
arXiv:1305.6226 (first posted May 27, 2013).\\
\textbf{Authors:} Jameson Cahill, Peter G. Casazza, Jesse Peterson, and
Lindsey Woodland.

The authors ask for examples of subspaces that
allow phase retrieval although the span of their associated projections is
not the span of any equally large family of rank-one projections
\cite{CCPW}.  This is Problem 5.9 on page 16 of the arXiv PDF (the source
labels it \texttt{P5}).

\begin{center}
  \includegraphics[width=0.96\linewidth]{figures/open_problem_crop.jpg}
\end{center}

\section*{Result}

This packet gives a \textbf{candidate full solution}: seven explicit
two-dimensional subspaces of \(\R^4\) allow phase retrieval, while the
seven-dimensional span of their orthogonal projections contains no nonzero
rank-one matrix.  The latter conclusion is stronger than the requirement in
the problem.

Here ``allows phase retrieval'' means that the map
\[
 x\in\R^4/\{\pm1\}\longmapsto
 (\lVert P_1x\rVert^2,\ldots,\lVert P_7x\rVert^2)
\]
is injective.  Equivalently, equality of all seven measurements forces
\(x=\pm y\).

\begin{theorem}\label{thm:main}
There are seven rank-two orthogonal projections \(P_1,\ldots,P_7\) on
\(\R^4\) such that
\begin{enumerate}
\item \(\{\operatorname{ran}P_i\}_{i=1}^7\) allows phase retrieval;
\item \(\spn\{P_1,\ldots,P_7\}\) contains no nonzero matrix of rank one.
\end{enumerate}
Consequently this span is not the span of any seven rank-one projections.
\end{theorem}

\section*{Idea of the proof}

Choose a three-dimensional subspace \(K\subset\Sym_4(\R)\) generated by
three pairwise anticommuting symmetric involutions.  The Clifford relations
make every nonzero member of \(K\) invertible.  Its orthogonal complement
\(S=K^\perp\) is then the span of seven explicit rank-two projections.
The phase-retrieval obstruction \(xx^{\mathsf T}-yy^{\mathsf T}\) has rank at
most two, so it cannot lie nontrivially in the measurement kernel \(K\).

The remaining requirement is more delicate: \(S\) must not be generated by
rank-one projections.  In fact it contains none.  Orthogonality of
\(zz^{\mathsf T}\) to the three Clifford generators produces the three
quadratic coordinates of the real Hopf map; their sum-of-squares identity
forces \(z=0\).

\section*{Construction}

Put
\[
 I=\begin{pmatrix}1&0\\0&1\end{pmatrix},\quad
 X=\begin{pmatrix}0&1\\1&0\end{pmatrix},\quad
 Z=\begin{pmatrix}1&0\\0&-1\end{pmatrix},\quad
 Y=\begin{pmatrix}0&-1\\1&0\end{pmatrix}.
\]
All tensor products below act on \(\R^2\otimes\R^2=\R^4\), in the standard
lexicographic basis.  Define
\[
 A_1=Z\otimes I,\qquad A_2=X\otimes I,\qquad A_3=Y\otimes Y
\]
and
\begin{align*}
 Q_1&=I\otimes X,& Q_2&=I\otimes Z,& Q_3&=X\otimes X,\\
 Q_4&=X\otimes Z,& Q_5&=Z\otimes X,& Q_6&=Z\otimes Z.
\end{align*}
Let \(I_4\) denote the identity on \(\R^4\), and set
\[
 P_1=\frac{I_4+Q_1}{2},\qquad
 P_2=\frac{I_4-Q_1}{2},\qquad
 P_{j+1}=\frac{I_4+Q_j}{2}\quad(2\leq j\leq6).
\]
Each \(Q_j\) is a symmetric traceless involution.  Hence every \(P_i\) is
an orthogonal projection of rank two.

For completeness, the seven requested subspaces can be read without tensor
notation from the following table.  Scaling the displayed generators is
irrelevant.
\[
\begin{array}{c@{\qquad}l}
\toprule
i&\operatorname{ran}P_i\\
\midrule
1&\spn\{(1,1,0,0),(0,0,1,1)\}\\
2&\spn\{(1,-1,0,0),(0,0,1,-1)\}\\
3&\spn\{(1,0,0,0),(0,0,1,0)\}\\
4&\spn\{(1,0,0,1),(0,1,1,0)\}\\
5&\spn\{(1,0,1,0),(0,1,0,-1)\}\\
6&\spn\{(1,1,0,0),(0,0,1,-1)\}\\
7&\spn\{(1,0,0,0),(0,0,0,1)\}\\
\bottomrule
\end{array}
\]

\section*{Proof}

Equip \(\Sym_4(\R)\) with the Hilbert--Schmidt inner product
\(\langle B,C\rangle=\operatorname{tr}(BC)\).  The ten matrices
\begin{equation}\label{eq:basis}
 I\otimes I,\ I\otimes X,\ I\otimes Z,\ X\otimes I,\ X\otimes X,
 \ X\otimes Z,\ Z\otimes I,\ Z\otimes X,\ Z\otimes Z,\ Y\otimes Y
\end{equation}
are nonzero, symmetric, and pairwise Hilbert--Schmidt orthogonal.  Since
\(\dim\Sym_4(\R)=10\), they form an orthogonal basis.

Let
\[
 K=\spn\{A_1,A_2,A_3\},\qquad S=K^\perp.
\]
Reading off the basis \eqref{eq:basis} gives
\begin{equation}\label{eq:S}
 S=\spn\{I_4,Q_1,\ldots,Q_6\}.
\end{equation}
Moreover,
\[
 I_4=P_1+P_2,\qquad Q_1=P_1-P_2,
 \qquad Q_j=2P_{j+1}-I_4\quad(2\leq j\leq6).
\]
Thus
\begin{equation}\label{eq:span}
 \spn\{P_1,\ldots,P_7\}=S.
\end{equation}

The matrices \(A_1,A_2,A_3\) are symmetric involutions which pairwise
anticommute.  It follows that, for all \(a,b,c\in\R\),
\begin{equation}\label{eq:clifford}
 (aA_1+bA_2+cA_3)^2=(a^2+b^2+c^2)I_4.
\end{equation}
In particular, every nonzero member of \(K\) is invertible and has rank four.

We now prove phase retrieval.  Suppose \(x,y\in\R^4\) satisfy
\(\lVert P_i x\rVert^2=\lVert P_i y\rVert^2\) for every \(i\).  With
\(D=xx^{\mathsf T}-yy^{\mathsf T}\), this says
\[
 \langle P_i,D\rangle=0\quad(1\leq i\leq7).
\]
By \eqref{eq:span}, \(D\in S^\perp=K\).  But \(\rank D\leq2\), whereas
\eqref{eq:clifford} says that every nonzero element of \(K\) has rank four.
Therefore \(D=0\).  Hence \(xx^{\mathsf T}=yy^{\mathsf T}\), which implies
\(x=\pm y\).  The seven subspaces allow phase retrieval.

It remains to prove the stronger rank-one exclusion.  If \(zz^{\mathsf T}\)
lies in \(S=K^\perp\), then
\(z^{\mathsf T}A_jz=0\) for \(j=1,2,3\).  Writing
\(z=(z_1,z_2,z_3,z_4)\), these three quadratic forms are
\begin{align*}
 q_1(z)&=z_1^2+z_2^2-z_3^2-z_4^2,\\
 q_2(z)&=2(z_1z_3+z_2z_4),\\
 q_3(z)&=2(z_1z_4-z_2z_3).
\end{align*}
They satisfy the elementary Hopf identity
\begin{equation}\label{eq:hopf}
 q_1(z)^2+q_2(z)^2+q_3(z)^2
   =(z_1^2+z_2^2+z_3^2+z_4^2)^2.
\end{equation}
Thus \(q_1(z)=q_2(z)=q_3(z)=0\) forces \(z=0\).  So \(S\) contains no
nonzero rank-one matrix.

Finally, if \(S\) were the span of seven rank-one projections
\(u_i u_i^{\mathsf T}\), every one of those nonzero generators would itself
belong to \(S\), contradicting the preceding paragraph.  This proves both
claims of Theorem~\ref{thm:main}. \qed

\section*{Machine-checkable audit}

The accompanying exact SymPy script checks symmetry, involution and
anticommutation; the determinant identity
\[
 \det(aA_1+bA_2+cA_3)=(a^2+b^2+c^2)^2;
\]
orthogonality and dimensions of \(K,S\); idempotence and rank of all seven
projections; equality of their span with \(S\); rank of the measurement map;
and the polynomial identity \eqref{eq:hopf}.  No floating-point computation
is used.

\section*{Novelty note}

A bounded search of the local corpus and arXiv/web indexes found later work
answering other projection-phase-retrieval questions, including generic
prescribed-rank constructions \cite{Edidin} and further structural results
\cite{CCJW}, but no source presenting this example or explicitly resolving
Problem 5.9.  This is a research packet, not a priority claim; specialist
bibliographic review remains necessary.

\section*{Limitations and review recommendation}

The result is an existence example in \(\R^4\), as requested; it does not
classify all such families or address the final complex minimality question
in the source paper.  The proof is finite-dimensional and exact, and every
matrix identity has been independently checked by symbolic computation.
Promotion for specialist review is recommended, with particular attention to
the span identification \eqref{eq:span} and the rank-one exclusion
\eqref{eq:hopf}.

\begin{thebibliography}{9}
\bibitem{CCPW}
J. Cahill, P. G. Casazza, J. Peterson, and L. Woodland,
\emph{Phase Retrieval by Projections}, arXiv:1305.6226.

\bibitem{Edidin}
D. Edidin, \emph{Projections and Phase Retrieval}, arXiv:1506.00674.

\bibitem{CCJW}
J. Cahill, P. G. Casazza, J. Jasper, and L. M. Woodland,
\emph{Phase Retrieval}, arXiv:1508.06687.
\end{thebibliography}

\end{document}
